% ============================================================================
% Section III: Problem Formulation and Controllers (draft for Overleaf)
% Needs in the preamble: amsmath, amssymb, graphicx, booktabs, xcolor,
% algorithm, algpseudocode. Figure file: fig_framework.pdf.
% Red [TBD ...] markers are numbers that do not exist yet; do not invent them.
% Labels kept from the previous draft: subsec:problem_formulation,
% subsec:belief_state, subsec:action_space, subsec:priced_objective,
% subsec:decision_rule, fig:system_flow, subsec:e1_results, subsec:e2_results.
% ============================================================================
\section{Problem Formulation and Controllers}
\label{sec:problem}

\subsection{Problem Formulation}
\label{subsec:problem_formulation}

A voice assistant answers one user turn at a time. While the user is still
speaking, a streaming speech recogniser produces partial transcripts
$x_1, x_2, \dots, x_K$, where $x_K$ is the transcript available when voice
activity detection (VAD) declares the end of the turn. The assistant may
retrieve documents before answering, or it may skip retrieval and answer from
the language model alone (closed-book). It therefore has to make two
decisions in every turn: \emph{whether} to retrieve, and \emph{when} to send
the retrieval request. A request sent while the user is still talking runs in
parallel with the remaining speech, so part of its latency is hidden. A
request sent too early, however, is built from an incomplete question and can
fetch the wrong documents.

Fig.~\ref{fig:system_flow} shows the pipeline. After each partial $x_k$ the
system updates two quantities: a belief $b_k$ that retrieval is needed for
this turn, and a prediction $\hat R_k$ of how much speech remains. A
controller reads them and either dispatches a retrieval, waits for the next
partial, or (at the end of the turn) skips retrieval. The retrieved passages,
if any, go to the generator (Mistral-7B), whose answer is spoken by a
text-to-speech engine. We measure the delay the user experiences as the time
to first audio,
\begin{equation}
\mathrm{TTFA}=t_{\text{first audio}}-t_{\text{user end}},
\label{eq:ttfa}
\end{equation}
where $t_{\text{user end}}$ is the VAD end of turn.

All controllers in this paper choose actions by maximising the same expected
utility, which rewards answer quality and charges for latency and compute
(Section~\ref{subsec:priced_objective}). No policy is learned by reinforcement
learning; each controller evaluates the objective directly. The controllers
differ only in how they organise the two decisions: a \emph{joint} controller
chooses whether and when in one decision, while a \emph{decomposed}
controller makes them in separate modules that do not share state
(Section~\ref{subsec:decision_rule}). The question this paper asks is when
that difference changes the outcome.

\begin{figure*}[t]
\centering
\includegraphics[width=\textwidth]{fig_framework.pdf}
\caption{Streaming pipeline and the decision structures compared. From each
ASR partial $x_k$ the system estimates the belief $b_k$ that retrieval is
needed and the remaining speech $\hat R_k$. (a) The joint controller decides
at every partial whether to dispatch now or wait, using both. (b) Decomposed
controllers fix the trigger without the belief and decide retrieval only at
that trigger. (c) Speculative controllers prefetch early and decide at the end
of the turn whether to use the result. All share one objective and are scored
by time to first audio (TTFA).}
\label{fig:system_flow}
\end{figure*}

\subsection{Belief State}
\label{subsec:belief_state}

The belief $b_k\in[0,1]$ is the controller's estimate of the probability that
the current turn needs retrieval. The two studies in this paper use
different belief signals.

\textit{Study 1 (v1).} The signal $s$ is the mean per-token entropy (in nats)
of a closed-book draft,
\begin{equation}
s=\frac{1}{N}\sum_{i=1}^{N}H\!\left(p_\theta(\cdot\mid y_{<i},q)\right),\qquad
b=\max\!\big(\beta,\ \sigma(w\,s+c)\big),
\label{eq:belief}
\end{equation}
averaged over five drafts (temperature 0.2, top-$p$ 0.9, 80 tokens). Here $q$
is the \emph{full} question, $\sigma$ is the logistic function, and
$\beta=0.05$ is a floor. The logistic map ($w=11.288$, $c=-3.633$) was fitted
on the 180 samples of Experiment~1 (36 items, five seeds) and was evaluated on
the same samples, so this calibration is in-sample. Maximum per-token entropy
ranked the items better than the mean
(Section~\ref{subsec:e1_results}). We kept the mean because its
distribution is smoother, which suits a logistic calibration. Because the v1 belief was computed once per item from the complete
question, it did not change with the dispatch time, even for policies that
dispatched before the question was finished.

\textit{Study 2 (v2).} Here an item counts as needing retrieval when
the model's closed-book answer is wrong in at least three of five sampled
drafts. In a real streaming system the belief can only use what
has been heard so far. We therefore recompute it on every partial $x_k$. The
signal is a self-assessment: we ask Mistral-7B whether the question can be
answered from general knowledge alone, without the organisation's own
information, and read the probability it assigns to ``No'' from its token
probabilities. This probe costs one forward pass instead of a full draft. The
probability is calibrated separately for each dataset on the development
split only, with a logistic model in the logit of the signal, the fraction of
the question heard so far, and their interaction. All test numbers use this
frozen calibration.

\subsection{Action Space and Dispatch Timing}
\label{subsec:action_space}

In v1 the controller chooses from four actions: skip, or retrieve with one of
three timing policies, $\mathcal T=\{\textsc{sentence\_end},
\textsc{immediate}, \textsc{fixed\_early\_proxy}\}$. Speech was simulated on
text at 150 words per minute, so a question of $n$ words lasts
$T=60n/150$~s. A prefix is \emph{query-ready} once it contains at least four
words, one of them a content word; let $t_{\text{ready}}$ be the first time
this holds.
\begin{itemize}
\item \textsc{sentence\_end} dispatches at $T$, with the full question.
\item \textsc{immediate} dispatches at the first partial after
$t_{\text{ready}}$, i.e.\ at $\min(\max(t_1,\ t_{\text{ready}}+\Delta),\ T)$,
where $t_1$ is the time of the first partial and $\Delta$ the partial
interval.
\item \textsc{fixed\_early\_proxy} dispatches at
$\min(\max(0.35T,\ t_{\text{ready}}),\ T)$, with $0.35T$ capped at $0.8\,t_1$.
\end{itemize}
Neither early policy predicts anything; \textsc{fixed\_early\_proxy} is a
fixed rule, not a forecast of the end of the turn.

In v2 the dispatch time is no longer one of a few fixed policies. At every
partial $k$ the controller may dispatch now or wait, and at the end of the
turn it either retrieves or skips. Deciding \emph{when} becomes an
optimal-stopping problem over the stream of partials.

We study two retrieval architectures. Under \textsc{commit}, a dispatched
retrieval is always used. Under \textsc{speculative}, a retrieval can be sent
early and then discarded at the end of the turn if the final belief says it
is not worth using; the compute it used is still paid. Both architectures are
run on the same items, seeds, latency profile and prices.

\subsection{Latency- and Compute-Priced Objective}
\label{subsec:priced_objective}

Retrieving when the belief is $b$ and the request is sent at time $t$ has
expected utility
\begin{equation}
U_t(b)=b\,g^{+}(1-\rho_t)-(1-b)\,g^{-}-\lambda_l\,e_t-\lambda_c\,(1+w_t),
\qquad U(\text{skip})=0 .
\label{eq:obj}
\end{equation}
Each term has a measured meaning:
\begin{itemize}
\item $g^{+}$ is the accuracy gained by retrieving on turns that need it, and
$g^{-}$ the accuracy lost by retrieving on turns that do not. Retrieval can
hurt: irrelevant passages can pull a correct closed-book answer off course.
\item $\rho_t$ is the fraction of $g^{+}$ lost because the query was built from
a partial transcript. It is measured as
$\rho(t)=1-\text{hit@3}(t)/\text{hit@3}(\text{end})$, and $\rho=0$ at the end
of the turn.
\item $e_t$ is the expected exposed latency in seconds, i.e.\ the part of
retrieval's cost the user still waits for (Eq.~\ref{eq:exposure}).
\item $w_t$ is the expected number of wasted retrievals caused by dispatching
early (for example a request that has to be re-sent).
\item $\lambda_l$ (accuracy per second) and $\lambda_c$ (accuracy per
retrieval) are prices that set how much latency and compute are worth
relative to accuracy.
\end{itemize}

The v1 objective was the special case $g^{+}=1$, $g^{-}=0$, $\rho_t=0$, which
sets the quality term to $b$. That form treats the probability that
retrieval is needed as if it were the size of the gain, and assumes retrieval
never hurts. Measured on the v1 scenarios, retrieval raised accuracy from
0.112 to 0.832 on turns that needed it ($g^{+}=0.720$, 95\% CI
[0.504, 0.888]) and lowered it from 0.704 to 0.112 on turns that did not
($g^{-}=0.592$, [0.408, 0.768]). At zero prices, Eq.~(\ref{eq:obj}) therefore
retrieves only when $b>b^{*}=g^{-}/(g^{+}+g^{-})=0.451$, not whenever
$b>0.05$ as the v1 form implied. In v2, $g^{\pm}$ are estimated on the
development split of each dataset. On d1\_open with recognised speech,
$\hat g^{+}=0.441$ and $\hat g^{-}=-0.112$: retrieval did not hurt items that
did not need it. On d2\_support, $\hat g^{+}=0.445$ and $\hat g^{-}=0.034$.

\textit{Exposed latency.} In v1, $e_t=L-o_t$, where $L$ is the retrieval
latency and $o_t$ is a constant overlap per timing policy. In v2 the overlap
depends on the turn:
\begin{equation}
e_k=\mathbb{E}_L\!\left[\max(0,\ L-W_k)\right]+\pi ,
\label{eq:exposure}
\end{equation}
where $W_k=\hat R_k+d$ is the predicted remaining speech plus the end-of-turn
detection delay $d$ ($W_K=0$). The term $\pi$ is the part of retrieval's cost
that no timing can hide, mainly the longer prompt the generator must read;
measured end to end with streaming TTS on the development split,
$\pi=371$~ms. We use three latency profiles: query rewriting followed by
FAISS search, by FAISS search with re-ranking, or by a remote retrieval API.
Their median retrieval latencies are 1{,}044, 1{,}043 and 1{,}241~ms. The
latency cost the controllers use is not the call duration but the end-to-end
TTFA increment, measured on the development split.

\textit{Waste rate.} In v1, $w_t$ was a placeholder (0, 0.02 or 0.08 per
policy) that we never measured: all 751 retrievals logged in v1 ran on text,
where no early request was ever invalidated, so the logs cannot estimate it.
In v2, $w$ is measured as the cancellation rate of early dispatches
\textcolor{red}{[TBD: measured rate]}, and we report how the results change
for $w\in\{0, 0.04, 0.08, 0.16\}$ \textcolor{red}{[TBD: sensitivity
analysis, Section~VII]}.

\textit{Prices.} Following v1, $\lambda_c=0.05$ throughout. In v1 we swept a
grid of $\lambda_l$ values. In v2 the primary test uses one pre-registered
price per dataset, the point of the grid at the 0.65 quantile of the
development-split end-of-turn belief: $\lambda_l=0.173$~s$^{-1}$ on
d1\_open and $0.216$~s$^{-1}$ on d2\_support.

\subsection{Controllers and Decision Rules}
\label{subsec:decision_rule}

Table~\ref{tab:controllers} lists the controllers. All priced controllers
use the same Eq.~(\ref{eq:obj}), the same belief calibration, the same
remaining-speech predictor, the same latency model and the same $g^{\pm}$.
They differ only in how the whether and when decisions are combined. All
thresholds and timing rules of the baselines are tuned on the development
split with an equal budget of evaluations.

\textbf{Joint, v1 (Ours).} The controller scores skip ($U=0$) and every
(retrieve, $t$) pair with Eq.~(\ref{eq:obj}) and takes the best; ties go to
skip.

\textbf{B9 (decomposed, full-latency pricing).} One module picks the timing
policy and another decides whether to retrieve. The two modules share no
state. Because the retrieval module does not know which timing will be used,
it prices retrieval at the full latency $L$ and never credits the overlap.
B9 can therefore skip a turn that would be worth retrieving with early
dispatch.

\textbf{B9+ (decomposed, expected-exposure pricing).} B9+ is still two separate
modules, but the retrieval module prices the expected exposure of the timing
its own timing module picks. This is the decomposed controller a careful
engineer would build, and it is the strongest decomposed rival under
\textsc{commit}. When the belief does not change with dispatch time and the
overlap is the same for every turn (the v1 setting; assumptions A1--A3 in
Section~IV), B9+ and Ours make the same decision: on v1 they agreed on all
11{,}300 decisions. Any advantage of joint control must therefore come from
the situations those assumptions exclude.

\textbf{Joint, v2 (Ours-v2).} Algorithm~\ref{alg:oursv2} treats dispatch as an
optimal-stopping problem. At each partial it compares the value of
dispatching now, $U_k(b_k)$, with the expected value of waiting for the next
partial. The next partial is assumed to carry new information with
probability $P(R>s)$, where $R$ is the predicted remaining speech and $s$ the
median spacing between partials on the development split. The belief
transitions are fitted on the development split. Ours-v2 has no tuned
parameter apart from its lookahead depth, which was fixed before the test.
Under \textsc{speculative} the same controller compares the value of having
the passages ready with the value of waiting.

\textbf{B10 (SPEC) and B11 (SPEC-$\theta$).} These are decomposed controllers
for the speculative architecture. B10 always prefetches at its trigger (the
first query-ready partial, or a tuned alternative) and, at the end of the
turn, uses the passages only if Eq.~(\ref{eq:obj}) evaluated at the final
belief $b_K$ is positive. B11 prefetches at the same trigger only if a
belief-only estimate of the prefetch's value, computed with fixed
development-split exposures, exceeds that of not prefetching; it then uses
B10's rule.

\textbf{Unpriced baselines.} B0 never retrieves. B4 and B6 are single-turn
proxies in the style of FLARE~\cite{jiang2023flare} and
DRAGIN~\cite{su2024dragin}: they retrieve at sentence end when one minus the
mean token confidence of the closed-book draft is at least 0.5 (B4), or when
its maximum token entropy is at least 1.0 (B6). B8 retrieves with a fixed
early dispatch when the maximum token entropy is at least 1.2. In v1 these
thresholds were fixed defaults and were not tuned; in v2 they are tuned on
the development split.

\begin{table}[t]
\caption{Controllers Compared. ``Priced'' means the controller reads
$(\lambda_l,\lambda_c)$; ``Joint'' means whether and when are chosen in one
decision.}
\label{tab:controllers}
\centering
\footnotesize
\setlength{\tabcolsep}{3pt}
\begin{tabular}{@{}l p{3.55cm} c c c@{}}
\toprule
ID & Retrieve rule / dispatch rule & Priced & Joint & Study\\
\midrule
B0 & never retrieve & -- & -- & v1, v2\\
B4 & $1-\bar p\ge0.5$ on closed-book draft (FLARE-style); sentence end & no & no & v1, v2\\
B6 & $s^{\max}\ge1.0$ (DRAGIN-style); sentence end & no & no & v1, v2\\
B8 & $s^{\max}\ge1.2$; fixed early dispatch & no & no & v1, v2\\
B9 & Eq.~(\ref{eq:obj}) with full latency $L$; timing chosen separately & yes & no & v1, v2\\
B9+ & Eq.~(\ref{eq:obj}) with the expected exposure of its own timing & yes & no & v1 (analysis), v2\\
B10 (SPEC) & always prefetch at trigger; use iff Eq.~(\ref{eq:obj}) $>0$ at end of turn & yes & no & v2\\
B11 (SPEC-$\theta$) & prefetch iff belief-only value (dev exposures) exceeds not prefetching; use rule of B10 & yes & no & v2\\
Ours & $\arg\max$ of Eq.~(\ref{eq:obj}) over skip and $(\text{retrieve},t)$ & yes & yes & v1\\
Ours-v2 & optimal stopping over partials (Alg.~\ref{alg:oursv2}) & yes & yes & v2\\
\bottomrule
\multicolumn{5}{@{}p{0.95\columnwidth}@{}}{\scriptsize v1 thresholds of B4, B6
and B8 were fixed defaults. v2 also runs B1, B5 and an oracle; all v2
thresholds are tuned on the development split with an equal budget.}\\
\end{tabular}
\end{table}

\begin{algorithm}[t]
\caption{Ours-v2: joint retrieve and dispatch decision for one turn
(\textsc{commit})}
\label{alg:oursv2}
\begin{algorithmic}[1]
\Require prices $(\lambda_l,\lambda_c)$; latency distribution of $L$;
belief-transition model fitted on dev; lookahead depth $n\ge1$
\For{$k=1,2,\dots$ \textbf{as} partial $k$ arrives at time $t_k$}
  \State observe $x_k$, $b_k$, $\hat R_k$ \Comment{only data available at $t_k$}
  \State $u\gets U_k(b_k)$ \Comment{Eq.~(\ref{eq:obj}) with $e_k$ from Eq.~(\ref{eq:exposure})}
  \If{$k=K$ (end of turn)}
    \State \Return \textsc{Retrieve} \textbf{if} $u>0$ \textbf{else} \textsc{Skip}
  \EndIf
  \State $v\gets\textsc{Wait}(b_k,\hat R_k,n)$ \Comment{expected value of waiting}
  \If{$u>0$ \textbf{and} $u\ge v$}
    \State \Return \textsc{Retrieve} \Comment{dispatch now}
  \EndIf
\EndFor
\end{algorithmic}
\end{algorithm}
